
What each objective was met with, and how the two research questions are answered, comes first (\S{}8.1). Then the contributions (\S{}8.2), each limitation paired with the experiment that would settle it (\S{}8.3), and a personal reflection on how the project actually went (\S{}8.4).

\section{Summary of the dissertation}

The aim was to work out whether spatial-relationship annotation for robot-acquired images can be automated, and whether the resulting labels are good enough to replace the human ones. Six objectives carried it.

\textbf{O1} is the pipeline, with 836 images annotated in about five minutes on a 6 GB GPU, no human deciding any label, and the dataset's formats written byte-compatibly. \textbf{O2} is the specification, with every threshold fitted on groups 0 to 5 alone. \textbf{O3} is the fidelity study, mean recall 0.85 and 0.74 on annotators no threshold saw, against three baselines, a vision-language model, ten ablations and audited precision. \textbf{O4} attributes every one of the 1,650 missed triplets to a cause, leaving roughly 7\% of them attributed to tool error (\S{}4.10, \S{}7.2). \textbf{O5} is the controlled experiment, 0.75 against 0.30 human and 0.36 self-trained. \textbf{O6} repeats it in a current SGG framework with a frozen detector and three seeds per arm, and carries it one link further to a planner. It is met in the sense that decides the answer, since the heavyweight test does not agree with the lightweight one, the disagreement is localised to the annotators whose labels Chapter 4 convicted of measured defects, and both readings are reported, not one of them selected. Section 1.2.2 fixed what would count as an answer before any result was reported. What follows is the verdict against those criteria.

\textbf{RQ1 is answered yes on five of seven predicates and qualified on two, against a criterion this dataset forced to be weaker than the one set.} The criterion was per-predicate recall on annotator groups no threshold ever saw, judged against two references fixed in advance. The trivial baselines are beaten by a wide margin, 0.74 held-out and 0.85 pooled against 0.14 for both random and majority, and the second reference, how well two human annotators would agree, is one this dataset cannot supply, for reasons Supplementary F.12 sets out.

So comparability rests on the baselines and the per-predicate audit. That is a weaker footing than the criterion intended. What the evidence establishes is that the tool reproduces the human record far above any baseline, and that the labels it adds beyond that record survive blind audit on five predicates. What it cannot establish is comparability with the human process itself, because no measurement of that process's own consistency exists to compare against. \textit{Answered yes} below is a verdict against the references that could be obtained, not against a human ceiling.

The second condition, that labels beyond the human record survive audit, holds for the lateral, proximity and depth-decided predicates at blind-audited precision 0.79--1.00. \textbf{It is not met for support.} At 0.535 [0.42, 0.65] about half the support labels the tool adds beyond the human record are wrong. That beats chance but falls short of a standard anything should be built on, so support answers RQ1 on recall and fails it on precision, and the per-predicate form of the question exists so that cannot be averaged away. For anything consuming the output, support is a candidate set, not a label set. The qualification is \texttt{in front of} and \texttt{behind} at 0.70/0.71 pooled, which is a disagreement about the words, not a failure of the criterion, with \S{}4.12 showing the tool reproducing its own verdict across viewpoints 0.958 of the time while two annotator groups applied the opposite convention.

\textbf{RQ2 is answered yes at the classifier, not met at the planner by the tool alone, and undecided at the benchmark.} The controlled classifier gives 0.75 against 0.30 on held-out \textit{human} gold, which is the rival source's own yardstick and the harder direction the criterion specified.

At the planner, \S{}1.2.2 fixed the comparator as the human arm. Against that arm's 25 of 25 the tool's relations alone clear 19, which falls short, and what draws level is the union with the vision-language source at 25 of 25, label-free but not the treatment RQ2 names. So the planner establishes that human annotation can be matched \textit{without a human in the loop}, not that the tool matches it.

At the benchmark, parity satisfies \textit{at least as well} as written. But a null result is not a demonstration of equivalence. Paired, the difference is -0.0006 with a 95\% interval of [-0.070, +0.069], so the honest word for 0.292 against 0.293 is \textit{undecided}, and since \S{}1.2.2 required agreement across all three for an unqualified yes, a tie is not agreement. What Chapter 5 measures as a two-and-a-half-fold advantage the ranked metric measures as parity. Section 6.3.1 localises the residual difference to the two test annotators with measured defects, and the third arm suits the argument least and is reported for it, with vision-language labels leading at 0.329 against 0.293 and 0.292 (\S{}6.3.2).

"Good enough to replace" is met at the classifier, unrefuted at the benchmark, and reached at the planner only once a second automatic source is added alongside. "Better" holds where ground truth means geometric consistency, not where it means annotation practice. The evidence carries this conditional, on this dataset: \textbf{automatic labels are better where ground truth means geometric consistency, and do not overtake human labels where it means annotation practice.} Robot planning needs the first. The planner is the one test that separates the sources at all, and it separates them in the human arm's favour until a second automatic source is added beside the tool's. On the ranked benchmark they never separate.

\section{Research contributions}

\textbf{For the source dataset and its authors.} An annotator that labels the existing images 20 times more densely in five minutes, in the dataset's own formats, with the operational definitions the annotation process never had. Six of its seven predicates are usable as labels, and support is a candidate set, for the reason \S{}4.8 gives. The fitted \texttt{near} threshold answers Wang et al.'s (2025) future-work request for "spatial thresholds for near", and two annotation defects are quantified for the first time: front/behind inverted in two of nine groups, and \texttt{near} supplied in quantity by only three, with fourfold variation in exhaustiveness.

\textbf{For work on scene-graph benchmarks.} Evidence, by dissociation, that ranked recall against sparse, guideline-free annotation partly measures agreement with annotator habits. Whichever model ranks better memorises which pairs annotators record, and the model that covers the relation types they never recorded is the one trained on consistent computed labels.

\textbf{For weak supervision.} A controlled three-way comparison, on the same features, model, split and seeds, showing that programmatic labels out-teach both scarce human labels and the standard remedy for them, in a feature space closely aligned with the labels' own generating mechanism, which is the scope \S{}5.5 sets and Chapter 6 was run to test. The mechanism is measured, not inferred, since self-training contributes roughly a thousand confident negative pseudo-labels for every positive one.

\textbf{Methodologically.} Calibration held out by \textit{annotator}, not by image. Sparse-gold protocol pairing recall with audited precision. Exhaustive loss attribution. One way to probe what annotators would score against one another without their ever having labelled the same images, together with the finding that on batches this unbalanced it yields an upper bound and no measurement (\S{}4.6). And a reliability check needing no labels at all, which came out of recovering the fact that the dataset was cut from a continuous capture (\S{}4.12). That last one is the most likely to transfer, since many robotics datasets are sequences presented as image sets, and there a predicate's agreement with itself across viewpoints separates a rule that is wrong from one that is merely uncertain.

\section{Limitations, and future research and development}

Each limitation below comes with the experiment that would settle it. Supplementary B specifies each of those experiments.

\textbf{The chain reaches the plan, not the robot.} The planner experiment (\S{}5.7) carries the chain one link further, from labels to the plan a robot would act on. But it runs on 25 occluder-selected scenes under a scorer that cannot see a false positive. So it is supporting evidence, not a closed link. What is missing outright is execution: no robot moved. The dataset ships no object meshes, so closing this needs a physics simulator (PyBullet or Isaac Sim) with primitive-shape proxies substituted for the detected boxes, then replaying the same 25 scenes' plans and scoring collision-free execution, not a text ordering. Execution scoring would also make the redundant-step cost of a hallucinated support relation visible, where now it is inferred from its precision.

\textbf{No precision figure is verdicted by a disinterested human.} Every one rests on the author's judgement and a vision-language model's, and while decoys in \S{}4.14 measure both judges and do not assume them, and the two disagree often enough that neither is echoing the other, two interested judges are not independence. Support is where this matters most, since it is also the predicate that fails. Settling it needs a panel with no stake in the tool, scoring a fresh draw from the shipped labels under the same blind instrument, and the sampling, rendering and scoring code are all in the repository.

\textbf{The benchmark is undecided, and the run that would decide it is written.} Three seeds bound the paired difference to [-0.070, +0.069], so \S{}6.7 reports parity without establishing it, while ten per arm would tighten that to about $\pm$0.020, and the notebook that would do it ships unrun for want of GPU-hours.

\textbf{One domain, one camera, six object classes.} What transfers is the calibration procedure, not the fitted constants, and the only out-of-domain evidence is qualitative (E.3).

\textbf{A class list stands in for geometry on \texttt{on}/\texttt{under}, with a known blind spot.} The support rule excludes \texttt{human} from either role, because mask contact alone cannot tell a person \textit{holding} an object from a surface \textit{supporting} one (0 of 2,466 gold support triplets involve a person either way). Ablation A10 tested the geometric alternative, a drop-fraction threshold. The two populations do not separate (Supplementary D.8). Scoped to the classes this dataset labels, the guard would be defeated by an object held by any unguarded class, whether a robot manipulator, a trolley or an animal, the same way a person defeats it. So every human-robot interaction scene this pipeline has not been run on carries that bound.

\textbf{Front/behind is bounded, but less by depth than the number suggests.} Section 4.9 bounds the engineering, since neither a larger depth model nor multi-frame geometry moves the pair, so the limit is monocular ambiguity in the scenes, not model capacity. A predicate reproducing its own verdict across viewpoints 0.958 of the time while recovering 0.70 of the human labels is applying a criterion the annotators did not share. So the intervention with the best expected return is a written annotation guideline ahead of a better network, which is an awkward conclusion for a computer-vision project and the one the evidence supports. What stays open is a calibrated stereo pair or an RGB-D capture. Wider surface detection was built, measured and declined (Supplementary D.4).

\textbf{Detection bounds full automation.} End-to-end recall with a zero-shot detector is 0.38 against the relation layer's 0.85 conditional on detection. Detection is the gap, not relations, and the source paper's own trained detector (0.93 mAP@50) would close most of it. Checking it needs only the released weights.

\textbf{Scale is shown without ground truth.} The pipeline was run over the 1,766 frames nobody has annotated: 562 keyframes after content-adaptive selection, 58 minutes, 185,242 triplets, a predicate distribution 0.032 in total variation from the annotated portion (E.5). Capacity and stability on unfamiliar input are therefore measured. Correctness on that portion is not, and cannot be without labels. Section 4.12 substitutes self-agreement for truth and should be read as the weaker thing it is.

\section{Personal reflections}

The most useful thing this project taught me was not to trust a number until I know how it was produced. That lesson arrived early. My first depth results were poor and every unit test passed, because the dataset's images are stored rotated 180 degrees behind an EXIF flag, so masks and depth were being read from an upside-down image while the boxes stayed upright. No automated check could have caught it, because every component was behaving as written. I found it by rendering an image with its boxes drawn on and looking at it, and I have inspected samples after every significant change since, which has caught two more. Lesson two was that "agreement with the humans" is not one target. I began by treating the human labels as ground truth and my disagreements as errors. Measuring them properly showed the three annotator behaviours Chapter 4 reports, and it reframed the question from "how close can I get to the humans" to "what do the humans mean, and where do they disagree". Almost every later design decision follows from that shift, including calibrating on some annotators and validating on others. Lesson three was about my own claims. On a single benchmark run I reported that my labels won against the one test annotator with clean conventions, by 0.011 over 73 images. Two more seeds showed them tied and I withdrew it (\S{}6.3.1). I had been careful in one place, fitting thresholds on some annotators and validating on others, and careless in another, treating one training run as a result. Recording the retraction was uncomfortable, and I think it was the right thing to do.

Given the time again I would front-load the experiments that answer the question the project actually asks. I spent considerable effort proving the labels are correct and train models well, and little on whether they help a robot decide what to do, which is the entire motivation. My planner experiment ran last and was the clearest single result, and run in week five it would have pointed the intervening work at supply of the support relation. I would also have had two people annotate the same fifty images in week two, since one afternoon would have produced outright the inter-annotator agreement figure I could only estimate indirectly. What I am most pleased with is not the accuracy figure but the two refinements I built, measured and declined. Both were plausible and took real work. The data said neither helped, and reporting that is where this started feeling like research.